% This LaTeX document needs to be compiled with XeLaTeX.
\documentclass[10pt]{jsarticle}
\usepackage{ascmac,amsmath,amssymb,amsthm,bm,physics,arydshln,mathcomp,wrapfig}
\newtheorem*{answer*}{解答}

\begin{document}
\textbf{「常微分方程式論 1 」 レポート課題}

\begin{flushright}
6122111 三森誠真
\end{flushright}


\begin{itembox}[l]{[1] 次の微分方程式を解け.}
\begin{equation*}
y^{2}+1+x^{2} y \frac{d y}{d x}=0 \tag{1}
\end{equation*}
  
\end{itembox}



\bigskip

\begin{answer*}
$x=0$ なら $y^{2}+1=0$ となり矛盾. ゆえに $x \neq 0$ である. (1) は変数分離形であるから,与式を変形して

$$
x^{2} y \frac{d y}{d x}=-\left(y^{2}+1\right), \quad \frac{y}{y^{2}+1} \frac{d y}{d x}=-\frac{1}{x^{2}}
$$

ここで両辺を $x$ について積分して


\begin{equation*}
\int \frac{y}{y^{2}+1} d y=-\int \frac{1}{x^{2}} d x+C_{0} \tag{*1}
\end{equation*}


このとき

$$
\begin{gathered}
\int \frac{y}{y^{2}+1} d y=\frac{1}{2} \int \frac{2 y}{y^{2}+1} d y=\frac{1}{2} \int \frac{\left(y^{2}+1\right)^{\prime}}{y^{2}+1} d y=\frac{1}{2} \log \left(y^{2}+1\right) \\
\int \frac{1}{x^{2}} d x=\frac{1}{-2+1} x^{-2+1}=-\frac{1}{x}
\end{gathered}
$$

であることに注意して, $(* 1)$ 式を書き換えると

$\frac{1}{2} \log \left(y^{2}+1\right)=\frac{1}{x}+C_{0}, \quad$ さらに変形して $\quad \log \left(y^{2}+1\right)=\frac{2}{x}+C_{1}, \quad$ with $C_{1}=2 C_{0}$

となる. ゆえに $y^{2}+1=\exp \left\{\frac{2}{x}+C_{1}\right\}=e^{C_{1}} e^{2 / x}=C_{2} e^{2 / x},\left(C_{2}=e^{C_{1}}\right)$ となるから, $y^{2}=C_{2} e^{2 / x}-1$ 

と書き直して, 最終的に

$$
y= \pm \sqrt{C e^{2 / x}-1}= \pm \sqrt{C \exp \left\{\frac{2}{x}\right\}-1} \quad \text { を得る. }
$$

\end{answer*}

\bigskip
\bigskip
\bigskip
\bigskip
\bigskip
\bigskip
\bigskip
\bigskip
\bigskip
\bigskip

\begin{itembox}[l]{$[2]$ 次の微分方程式を解け.}
 \begin{equation*}
y^{2}-3 x^{2}+2 x y \frac{d y}{d x}=0 \tag{2}
\end{equation*}
 
\end{itembox}

\begin{answer*}
\end{answer*}

与式 (2) において $x=0$ のケースを考えると, $y^{2}=0$ より $y=0$ (ゼロ解) が従う.ゼ 口解は 自明なので $x \neq 0$ としてよい.

ゆえに (2) の両辺に $\times \frac{1}{x^{2}}$ をかけて変形して


\begin{equation*}
\left(\frac{y}{x}\right)^{2}-3+2\left(\frac{y}{x}\right) \frac{d y}{d x}=0 \tag{*0}
\end{equation*}


を得る．そこで $(* 0)$ 式をさらに変形して


\begin{equation*}
2\left(\frac{y}{x}\right) \frac{d y}{d x}=3-\left(\frac{y}{x}\right)^{2}, \quad \frac{d y}{d x}=\frac{3-\left(\frac{y}{x}\right)^{2}}{2\left(\frac{y}{x}\right)}=F\left(\frac{y}{x}\right) \text { 形となる. } \tag{*1}
\end{equation*}

ここで

\begin{equation*}
\frac{y}{x}=u \quad \text { とおいて }, \quad y=x u, \quad y^{\prime}=u+x u^{\prime} \tag{*2}
\end{equation*}


であるから,これにより方程式 $(* 0)$ を書き直して

$$
u^{2}-3+2 u \frac{d y}{d x}=0, \quad u^{2}-3+2 u\left(u+x u^{\prime}\right)=0
$$

を得る. さらにこれを整理して

$$
\begin{gathered}
2 x u \cdot u^{\prime}=3\left(1-u^{2}\right) \quad \text { となるので，これは } \\
\frac{u}{1-u^{2}} \cdot \frac{d u}{d x}=\frac{3}{2} \cdot \frac{1}{x} \quad \text { と変形して,変数分離形である. } \\
-\frac{1}{2} \int \frac{-2 u}{1-u^{2}} d u=\frac{3}{2} \int \frac{1}{x} d x+C_{0}=\frac{3}{2} \log |x|+C_{0}
\end{gathered}
$$

したがって, $-\frac{1}{2} \log \left|1-u^{2}\right|=\frac{3}{2} \log |x|+C_{0}$ を得る.ここで両辺を 2 倍して

$-\log \left|1-u^{2}\right|=3 \log |x|+2 C_{0}=\log |x|^{3}+\log e^{2 C_{0}}=\log C_{1} x^{3} \quad$ with $\quad C_{1}= \pm e^{2 C_{0}}$ $\log \left|1-u^{2}\right|=-\log C_{1} x^{3}=\log \frac{1}{C_{1} x^{3}}, \quad 1-u^{2}=\frac{1}{C_{2} x^{3}} \quad$ with $\quad C_{2}= \pm C_{1}$

ゆえに


\begin{equation*}
u^{2}=1-\frac{1}{C_{2} x^{3}} \text { を得る. } \tag{*3}
\end{equation*}


ここで $u=y / x$ を代入して $, x, y$ の式に戻して

$$
\left(\frac{y}{x}\right)^{2}=\frac{y^{2}}{x^{2}}=1-\frac{1}{C_{2} x^{3}} \quad y^{2}=x^{2}-\frac{1}{C_{2} x}=x^{2}+\frac{C}{x} \quad \text { with } \quad C=-\frac{1}{C_{2}}
$$

かくして方程式 $(2)$ の解

$$
y^{2}=x^{2}+\frac{C}{x}, \quad \text { or } \quad y= \pm \sqrt{\frac{x^{3}+C}{x}}= \pm x \sqrt{1+\frac{C}{x^{3}}}
$$

を得る。

\bigskip
\bigskip
\bigskip
\bigskip
\bigskip
\bigskip
\bigskip
\bigskip

\begin{itembox}[l]{$[3]$ 次の微分方程式を解け.}
\begin{equation*}
x-2 y+5+(2 x-y+4) \frac{d y}{d x}=0 \tag{3}
\end{equation*} 
(ヒント) $X=x+1, Y=y-2$ と変数変換して, $X, Y$ の方程式に書き換え, 同次形の 微分方程式


\begin{equation*}
\frac{d Y}{d X}=F\left(\frac{Y}{X}\right) \tag{4}
\end{equation*}


を導くことによって解け.
\end{itembox}


\bigskip

\begin{answer*}
\end{answer*}

元の方程式 $(3)$ を書き換えて


\begin{equation*}
(2 x-y+4) \frac{d y}{d x}=-(x-2 y+5), \quad \frac{d y}{d x}=-\frac{x-2 y+5}{2 x-y+4} \tag{*1}
\end{equation*}


を得る.ここで

\[
\left\{\begin{array}{l}
x-2 y+5=0  \tag{*2}\\
2 x-y+4=0
\end{array}\right.
\]

を連立させて解いて, $x=-1, y=2$ を得る.したがって，

\[
\left\{\begin{array}{l}
X=x+1  \tag{*3}\\
Y=y-2
\end{array}\right.
\]

と変数変換して, 与えられた微分方程式 $(3)$ 或いは $(* 1)$ を書き直して, $(X, Y)$ に関する微分方程式を導く. 実際, $X-2 Y=x-2 y+5,2 X-Y=2 x-y+4$ であること,および $d Y=d y$ かつ $d X=d x$ であることに注意すれば, 新しい $(X, Y)$ に関する方程式


\begin{equation*}
\frac{d Y}{d X}=-\frac{X-2 Y}{2 X-Y} \tag{*4}
\end{equation*}


を得る.この $(* 4)$ 式を少し変形すると


\begin{equation*}
\frac{d Y}{d X}=-\frac{X-2 Y}{2 X-Y}=-\frac{1-2 \frac{Y}{X}}{2-\frac{Y}{X}} \tag{*5}
\end{equation*}

この微分方程式 $(* 5)$ は同次形であるから, $Z=Y / X$ とお くことで変数分離形に帰着できる.

$$
Y=X Z, \quad \frac{d Y}{d X}=Z+X \frac{d Z}{d X}
$$

に注意して, $(* 5)$ から直ちに $Z$ に関する微分方程式


\begin{equation*}
Z+X \frac{d Z}{d X}=-\frac{1-2 Z}{2-Z} \tag{*6}
\end{equation*}


を得る、これを整理して $X Z^{\prime}=\left(Z^{2}-1\right) /(2-Z)$ となるから, 変数分離形である.したがって, 変形して


\begin{equation*}
\frac{2-Z}{Z^{2}-1} \cdot \frac{d Z}{d X}=\frac{1}{X}, \quad \text { 両辺を } X \text { について積分すると } \quad \int \frac{2-Z}{Z^{2}-1} d Z=\int \frac{1}{X} d X+C_{0} \tag{*7}
\end{equation*}


を得る.一方，


\begin{equation*}
I_{1}:=\int \frac{2-Z}{Z^{2}-1} d Z=2 \int \frac{d Z}{Z^{2}-1}-\int \frac{Z}{Z^{2}-1} d Z=: 2 J_{1}-J_{2} \quad \text { とおく. } \tag{*8}
\end{equation*}


$$
\begin{aligned}
& J_{1}=\int \frac{d Z}{Z^{2}-1}=\int \frac{d Z}{(Z+1)(Z-1)}=-\frac{1}{2} \int \frac{d Z}{Z+1}+\frac{1}{2} \int \frac{d Z}{Z-1}=\frac{1}{2} \log \left|\frac{Z-1}{Z+1}\right| \\
& J_{2}=\int \frac{Z}{Z^{2}-1} d Z=\frac{1}{2} \int \frac{2 Z}{Z^{2}-1} d Z=\frac{1}{2} \int \frac{\left(Z^{2}-1\right)^{\prime}}{Z^{2}-1} d Z=\frac{1}{2} \log \left|Z^{2}-1\right| \\
& I_{2}:=\int \frac{1}{X} d X+C_{0}=\log |X|+C_{0}
\end{aligned}
$$

であることに注意すれば, $(* 7)$ および $(* 8)$ から関係式


\begin{equation*}
2 \times \frac{1}{2} \log \left|\frac{Z-1}{Z+1}\right|-\frac{1}{2} \log \left|Z^{2}-1\right|=\log |X|+C_{0} \tag{*9}
\end{equation*}


を得る. これを整理してまとめると

$$
\log \frac{(Z-1)^{2}}{\left(Z^{2}-1\right)(Z+1)^{2}}=\log C_{1} X^{2}, \quad C_{1}=e^{2 C_{0}}
$$

さらに対数関数の一意性より


\begin{equation*}
\frac{Z-1}{(Z+1)^{3}}=C X^{2} \tag{*10}
\end{equation*}


を得る.ここで $Z=Y / X$ を用いて, 元の変数に戻して

$$
\frac{\frac{Y}{X}-1}{\left(\frac{Y}{X}+1\right)^{3}}=C X^{2}, \quad \Longrightarrow \quad \frac{Y-X}{(Y+X)^{3}}=C, \quad\left\{\begin{array}{l}
X=x+1 \\
Y=y-2
\end{array}\right. \text { を代入して }
$$

さらに $Y-X=y-x-3$, かつ $Y+X=y+x-1$ に注意して, 元の方程式 (3) の解

$$
\frac{y-x-3}{(y+x-1)^{3}}=c \quad \text { を得る. }
$$

一方, 符号を変えて, 任意定数に $(-1)$ を繰り込むと

$$
\frac{X-Y}{(X+Y)^{3}}=C \quad \text { とできるから }, \quad \frac{x-y+3}{(x+y-1)^{3}}=c
$$

としてもよい.

\bigskip
\bigskip
\bigskip
\bigskip
\bigskip
\bigskip
\bigskip
\bigskip
\bigskip
\bigskip
\bigskip
\bigskip
\bigskip
\bigskip
\bigskip
\bigskip

\begin{itembox}[l]{[4] 次の微分方程式を解け.}
\begin{equation*}
x \frac{d y}{d x}+(1-x) y=1 \tag{5}
\end{equation*}
\end{itembox}

\bigskip

\begin{answer*}
\end{answer*}

(5) は 1 階線形微分方程式であるから, 定数変化法により解く.

まず (5) の斉次微分方程式: $x \frac{d z}{d x}+(1-x) z=0 \quad \cdots \cdots(* 1)$ を解く.

移行して $x \frac{d z}{d x}=(x-1) z, \quad$ これは 変数分離形であるから, 変形して $\frac{1}{z} \frac{d z}{d x}=1-\frac{1}{x}$ さらに両辺を $x$ について積分して

$$
\int \frac{1}{z} d z=\int\left(1-\frac{1}{x}\right) d x+C_{0}=\int d x-\int \frac{1}{x} d x+C_{0}
$$

不定積分を実行して $\quad \log |z|=x-\log |x|+C_{0}=\log e^{x}-\log |x|+\log e^{C_{0}}=\log \frac{C_{1} e^{x}}{|x|}$ ここで $C_{1}=e^{C_{0}}$ と置いた. 対数関数の 1 対 1 対応性から

$$
z=C \frac{e^{x}}{x} \cdots \cdots(* 2) \quad \text { を得る. } \quad\left(C= \pm C_{1}\right)
$$

次に元々の非斉次方程式 (5) の一般解を求める. $(* 2)$ に対して, 定数変化法を 用いて (5) の解を


\begin{equation*}
y \equiv y(x)=f(x) \cdot \frac{e^{x}}{x} \tag{*3}
\end{equation*}


として求める. $(* 3)$ を微分して


\begin{equation*}
y^{\prime}(x)=f^{\prime}(x) \cdot \frac{e^{x}}{x}+f(x)\left\{\frac{e^{x}}{x}\right\}^{\prime}=f^{\prime}(x) \cdot \frac{e^{x}}{x}+f(x) \cdot \frac{e^{x}}{x}-f(x) \cdot \frac{e^{x}}{x^{2}} \tag{*4}
\end{equation*}


を得る. よって元の方程式 (5) に代入して

$$
x\left\{f^{\prime}(x) \cdot \frac{e^{x}}{x}+f(x) \cdot \frac{e^{x}}{x}-f(x) \cdot \frac{e^{x}}{x^{2}}\right\}+(1-x) f(x) \cdot \frac{e^{x}}{x}=1, \quad f^{\prime}(x) e^{x}=1
$$

を得るから $, f^{\prime}(x)=e^{-x}, \cdots \cdots(* 5)$ を積分して


\begin{equation*}
f(x)=\int e^{-x} d x+C_{2}=-e^{-x}+C_{2} \tag{*6}
\end{equation*}


を得る. ゆえに $(* 6)$ を $(* 3)$ に代入して


\begin{equation*}
y=f(x) \cdot \frac{e^{x}}{x}=\left(-e^{-x}+C_{2}\right) \frac{e^{x}}{x}=-\frac{1}{x}+C_{2} \frac{e^{x}}{x} \tag{*7}
\end{equation*}


となる.したがって, 求める非斉次方程式 (5) の一般解は


\begin{equation*}
y(x)=-\frac{1}{x}+C \frac{e^{x}}{x} \tag{*8}
\end{equation*}


で与えられる.

\bigskip

\begin{itembox}[l]{$[5]$ 次の微分方程式を解け.}
\begin{equation*}
\frac{d y}{d x}+2 x y=2 x^{3} y^{3} \tag{6}
\end{equation*}  
\end{itembox}

\bigskip

\begin{answer*}
\end{answer*}

(6) 式の両辺に $\times \frac{1}{y^{3}}$ を掛けて,


\begin{equation*}
\frac{y^{\prime}}{y^{3}}+\frac{2 x}{y^{2}}=2 x^{3} \tag{*1}
\end{equation*}


となる.ここで $z=1 / y^{2}$ と置いて, 方程式 $(* 1)$ を $z$ に関する微分方程式に書き直す.

$$
y^{2}=\frac{1}{z}=z^{-1}, \text { 微分して } 2 y \frac{d y}{d x}=(-1) z^{-2} \frac{d z}{d x}
$$

直ちに $y^{\prime}=\frac{1}{2 y} \cdot(-1) \frac{1}{z^{2}} z^{\prime}, \quad$ さらに整理して,$\quad \Longrightarrow \quad \frac{y^{\prime}}{y^{3}}=-\frac{1}{2} z^{\prime}$

を得る. したがって,これらを $(* 1)$ に代入して


\begin{equation*}
-\frac{1}{2} z^{\prime}+2 x z=2 x^{3}, \quad \Longrightarrow \quad z^{\prime}-4 x z=-4 x^{3} \tag{*3}
\end{equation*}


を得る.この $(* 3)$ 式は $z$ に関する 1 階線形微分方程式に他ならない. したがって定数変化法により解くことができる.

まず次の斉次方程式を解く. $w^{\prime}-4 x w=0$. ここで $w^{\prime}=4 x w$ は変数分離形 であるから, 変形して

$$
\frac{1}{w} \frac{d w}{d x}=4 x, \quad \text { 両辺を } x \text { について積分すると } \quad \int \frac{1}{w} d w=4 \int x d x+C_{0}
$$

であるから,

$$
\log |w|=4 \cdot \frac{x^{2}}{2}+C_{0}=2 x^{2}+C_{0}=\log \exp \left\{2 x^{2}+C_{0}\right\}=\log e^{C_{0}} e^{2 x^{2}}=\log C_{1} \cdot e^{2 x^{2}}
$$

ここで $C_{1}:=e^{C_{0}}$ と置いた. これより直ちに斉次方程式の解


\begin{equation*}
w(x)= \pm C_{1} e^{2 x^{2}}=C_{2} e^{2 x^{2}}, \quad C_{2}= \pm C_{1} \tag{*4}
\end{equation*}


が得られる.

次に定数変化法を用いて非斉次方程式の一般解を求める. 上記の結果 $(* 4)$ を 受けて, 非斉次方程式 $(* 3)$ の解を


\begin{equation*}
z(x)=f(x) e^{2 x^{2}} \tag{*5}
\end{equation*}


とする. 微分して $z^{\prime}=f^{\prime} e^{2 x^{2}}+f \cdot e^{2 x^{2}} \times(4 x)=f^{\prime} e^{2 x^{2}}+4 x z$. この式を上の $(* 3)$ 式に 代入すれば, 直ちに $f^{\prime} e^{2 x^{2}}=-4 x^{3}$ を得る. 結局, $f^{\prime}=-4 x^{3} e^{-2 x^{2}}$ であるから, 積分して $f(x)$ を求める. 実際,

$$
f(x)=-4 \int x^{3} e^{-2 x^{2}} d x+C
$$

であるから, 求める非斉次方程式 $(* 3)$ の解は


\begin{equation*}
z(x)=e^{2 x^{2}}\left(-4 \int x^{3} e^{-2 x^{2}} d x+C\right) \tag{*6}
\end{equation*}


で与えられる. あとは上述の $(* 6)$ 内の不定積分 $I=\int x^{3} e^{-2 x^{2}} d x$ を具体的に求めさえすればよい.

$$
\left(e^{-2 x^{2}}\right)^{\prime}=-4 x \cdot e^{-2 x^{2}}
$$

に注意して, 部分積分法を用いると


\begin{align*}
I & =\int x^{3} e^{-2 x^{2}} d x=\int x^{2} \cdot x e^{-2 x^{2}} d x=\frac{-1}{4} \int x^{2} \cdot(-4 x) e^{-2 x^{2}} d x \\
& =\frac{-1}{4} \int x^{2} \cdot\left(e^{-2 x^{2}}\right)^{\prime} d x=\frac{-1}{4}\left\{x^{2} e^{-2 x^{2}}-\int 2 x e^{-2 x^{2}} d x\right\} \\
& =\frac{-1}{4} x^{2} e^{-2 x^{2}}+\frac{1}{2} \int x e^{-2 x^{2}} d x  \tag{*7}\\
J & :=\int x e^{-2 x^{2}} d x=\frac{-1}{4} \int(-4) x \cdot e^{-2 x^{2}} d x=\frac{-1}{4} \int\left(e^{-2 x^{2}}\right)^{\prime} d x \\
& =\frac{-1}{4} e^{-2 x^{2}}+C \tag{*8}
\end{align*}


ゆえに, $(* 7)$ と $(* 8)$ の結果を合わせて

$$
I=\int x^{3} e^{-2 x^{2}} d x=\frac{-1}{4} x^{2} e^{-2 x^{2}}+\frac{1}{2} \times \frac{-1}{4} e^{-2 x^{2}}+C
$$

であるから, $(* 6)$ 式に代入して, 求める非斉次方程式 $(* 3)$ の一般解は


\begin{equation*}
z=e^{2 x^{2}}\left(x^{2} e^{-2 x^{2}}+\frac{1}{2} e^{-2 x^{2}}+C\right)=x^{2}+\frac{1}{2}+C e^{2 x^{2}} \tag{*9}
\end{equation*}


を得る。

元々 Bernoulli 型微分方程式 (6) を考えていたので式 $z=1 / y^{2}$ を考慮して


\begin{equation*}
y^{2}=\frac{1}{z}=\frac{1}{x^{2}+\frac{1}{2}+C e^{2 x^{2}}} \tag{*10}
\end{equation*}


である. あるいは, 求める一般解は


\begin{equation*}
y(x)=\frac{1}{ \pm \sqrt{x^{2}+\frac{1}{2}+C e^{2 x^{2}}}} \tag{*11}
\end{equation*}


となる.

\bigskip
\bigskip
\bigskip
\bigskip
\bigskip
\bigskip
\bigskip
\bigskip
\bigskip
\bigskip
\bigskip
\bigskip
\bigskip
\bigskip

\begin{itembox}[l]{[6]}
次の微分方程式は $\frac{a}{x}$ の形の特殊解をもつという.このとき,この方程式の一般解を求めよ.

\begin{equation*}
\frac{d y}{d x}+y^{2}=\frac{2}{x^{2}} \tag{7}
\end{equation*}
\end{itembox}





\bigskip

\begin{answer*}
\end{answer*}

特殊解 $y=a / x$ は微分方程式 (7) をみたすから, パラメータ $a$ を決めよう.


\begin{equation*}
y^{\prime}=\left(\frac{a}{x}\right)^{\prime}=\left(a x^{-1}\right)^{\prime}=-a \cdot x^{-2}=-\frac{a}{x^{2}} \tag{*1}
\end{equation*}


この $(* 1)$ を方程式 $(7)$ に代入して

$$
-\frac{a}{x^{2}}+\left(\frac{a}{x}\right)^{2}=\frac{2}{x^{2}} \quad \Longrightarrow \quad-\frac{a}{x^{2}}+\frac{a^{2}}{x^{2}}=\frac{-a+a^{2}}{x^{2}}=\frac{2}{x^{2}}
$$

これより $a$ は関係式 $a^{2}-a=2$ をみたすことがわかる. 方程式 $a^{2}-a-2=0$ を解い $\tau,(a+1)(a-2)=0$, したがって, $a=-1$ あるいは $a=2$ である. ここでは,


\begin{equation*}
a=2 \tag{*2}
\end{equation*}


を採用して話を進めることにする. (7) はリッカチ型微分方程式であるから, 特殊解が既知の元での解法の処方箋に従って


\begin{equation*}
y=\frac{2}{x}+u \tag{*3}
\end{equation*}


と置いて, 未知関数を $u$ に変換して, $u$ がみたすべき方程式を導き出す. そこで $(* 3)$ を元 の方程式 $(7)$ に代入して


\begin{align*}
& y^{\prime}=\left(\frac{2}{x}\right)^{\prime}+u^{\prime}=-\frac{2}{x^{2}}+u^{\prime}, \quad-\frac{2}{x^{2}}+u^{\prime}+\left(\frac{2}{x}+u\right)^{2}=\frac{2}{x^{2}} \\
& \left(\frac{2}{x}\right)^{\prime}+u^{\prime}+\left(\frac{2}{x}\right)^{2}+4 \frac{u}{x}+u^{2}=\frac{2}{x^{2}} \tag{*4}
\end{align*}


となる. ここで $y=2 / x$ が特殊解であり, 特殊解 $y=2 / x$ は元の微分方程式 (7) $y^{\prime}+y^{2}=2 / x^{2}$ をみたすから, 上の式 $(* 4)$ で

$$
\left(\frac{2}{x}\right)^{\prime}+\left(\frac{2}{x}\right)^{2}=\frac{2}{x^{2}}
$$

なる関係をみたし, 結局,$(* 4)$ 式から


\begin{equation*}
u^{\prime}+4 \frac{u}{x}+u^{2}=0 \tag{*5}
\end{equation*}


という $u$ に関する新しい微分方程式が導き出される。これはベルヌーイ (Bernoulli) 型微分方程式に他ならない.したがって, 適当な変換により1階線形微分方程式に帰着できる.

実際, 方程式 $(* 5)$ の両辺に $\times \frac{1}{u^{2}}$ を掛けて変形すると

$$
\frac{u^{\prime}}{u^{2}}+\frac{4}{x} \cdot \frac{1}{u}+1=0
$$

を得る.ここで $(1 / u)^{\prime}=\left(u^{-1}\right)^{\prime}=(d / d x)\left(u^{-1}\right)=(-1) u^{-2} \cdot(d u / d x)$ であることに注意 して, 少し見方を変えて, 上の式から微分関係式

$$
-\left(\frac{1}{u}\right)^{\prime}+\frac{4}{x}\left(\frac{1}{u}\right)=-1
$$

が成立していることがわかる. すなわち，


\begin{equation*}
\left(\frac{1}{u}\right)^{\prime}-\frac{4}{x}\left(\frac{1}{u}\right)=1 \tag{*6}
\end{equation*}


を得る. つまり, 変換 $z=1 / u, \cdots \cdots(* 7)$ を考えれば,


\begin{equation*}
z^{\prime}-\frac{4}{x} z=1 \tag{*8}
\end{equation*}


という 1 階線形微分方程式に帰着された.これは定数変化法により解くことができる.

$\quad((* 8)$ の斉次方程式を解くこと)

$$斉次
w^{\prime}-\frac{4}{x} w=0, \quad \Longrightarrow \quad \frac{d w}{d x}=\frac{4}{x} w
$$

これは変数分離形であるから, 変形して, 両辺を $x$ について積分して

$$
\begin{aligned}
& \frac{1}{w} \frac{d w}{d x}=\frac{4}{x}, \quad \int \frac{d w}{w}=\int \frac{4}{x} d x+C_{0} \\
& \log |w|=4 \log |x|+C_{0}=\log x^{4}+\log e^{C_{0}}=\log C_{1} x^{4}, \quad\left(C_{1}=e^{C_{0}}\right) \\
& \Longrightarrow \quad w \equiv w(x)= \pm C_{1} x^{4}=C_{2} x^{4} . \quad \text { を得る. }
\end{aligned}
$$

(定数変化法により非斉次方程式の一般解を求めること）

非斉次方程式 $(* 8) z^{\prime}-(4 / x) z=1$ の解を定数変化法により


\begin{equation*}
z \equiv z(x)=f(x) x^{4} \tag{*9}
\end{equation*}


として, この $f(x)$ を定める. 実際, $z^{\prime}=f^{\prime} x^{4}+4 x^{3} \cdot f$ であるから, これらを元の方程式 $(* 8)$ に代入して

$$
f^{\prime} x^{4}+4 x^{3} \cdot f-\frac{4}{x} \cdot f x^{4}=1, \quad \text { キャンセルにより } \quad f^{\prime} x^{4}=1
$$

これより $f^{\prime}(x)=1 / x^{4}, \cdots \cdots(* 10)$ を得る.この $(* 10)$ を $x$ について積分して


\begin{align*}
f(x) & =\int \frac{1}{x^{4}} d x+C=\int x^{-4} d x+C=\frac{1}{-4+1} x^{-4+1}+C=\frac{1}{-3} x^{-3}+C=\frac{-1}{3 x^{3}}+C \\
z & =x^{4}\left(\frac{-1}{3 x^{3}}+C\right)=-\frac{x}{3}+C x^{4}=C x^{4}-\frac{x}{3} \quad(* 11) \tag{*11}
\end{align*}


この $(* 11)$ は非斉次方程式 $(* 8)$ の解を与える.


\begin{equation*}
z=\frac{1}{u} \tag{*7}
\end{equation*}


と置いたので, 上の $(* 11)$ の結果より


\begin{equation*}
u=\frac{1}{z}=\frac{1}{C x^{4}-\frac{x}{3}} \tag{*12}
\end{equation*}


である. ゆえに求めるリッカチ (Riccati) 型微分方程式 (7)

$$
\frac{d y}{d x}+y^{2}=\frac{2}{x^{2}}
$$

の一般解は

$$
\begin{aligned}
y & =\frac{2}{x}+u=\frac{2}{x}+\frac{1}{C x^{4}-\frac{x}{3}}=\frac{2 C x^{4}-\frac{2}{3} x+x}{x\left(C x^{4}-\frac{x}{3}\right)} \\
& =\frac{2 C x^{4}+\frac{x}{3}}{C x^{5}-\frac{x^{2}}{3}}=\frac{6 C x^{4}+x}{3 C x^{5}-x^{2}}=\frac{6 C x^{3}+1}{3 C x^{4}-x}=\frac{2 c x^{3}+1}{c x^{4}-x}
\end{aligned}
$$

ただし,ここで $c=3 C$ と置いた. かくして求めるリッカチ型方程式 (7) の解は


\begin{equation*}
y=\frac{2 c x^{3}+1}{c x^{4}-x} \tag{*13}
\end{equation*}


で与えられる.

\bigskip
\bigskip
\bigskip
\bigskip
\bigskip
\bigskip
\bigskip

\begin{itembox}[l]{$[7]$ 次の微分方程式の一般解を求めよ.}
\begin{align*}
& \text { (1) } y^{\prime \prime}+y^{\prime}-2 y=0  \tag{8}\\
& \text { (2) } y^{\prime \prime}+2 y^{\prime}+2 y=0  \tag{9}\\
& \text { (3) } y^{\prime \prime}+2 y^{\prime}+y=0 \tag{10}
\end{align*}  
\end{itembox}

\bigskip
\bigskip

\begin{answer*}
\end{answer*}

(1) $y^{\prime \prime}+y^{\prime}-2 y=0$ で $y=e^{\alpha x}$ と置くと, $y^{\prime}=\alpha e^{\alpha x}, y^{\prime \prime}=\alpha^{2} e^{\alpha x}$ より上方程式に代入して

$$
\left(\alpha^{2}+\alpha-2\right) e^{\alpha x}=0, \quad e^{\alpha x}>0 \text { より } \quad \alpha^{2}+\alpha-2=0 \text { (特性方程式) }
$$

を得る.この特性方程式を解いて, $(\alpha-1)(\alpha+2)=0$ より直ちに $\alpha=1$ あるいは $\alpha=-2$ を得る.このとき, $e^{x}$ と $e^{-2 x}$ は互いに独立な解であるから, 方程式 (1) の一般解はこれ らの 1 次結合で与えられる. すなわち,

$$
y=c_{1} e^{x}+c_{2} e^{-2 x}, \quad\left(c_{1}, c_{2}: \text { 任意定数 }\right)
$$

である.

(2) $y^{\prime \prime}+2 y^{\prime}+2 y=0$ で $y=e^{\alpha x}$ と置いて, $y^{\prime}=\alpha e^{\alpha x}, y^{\prime \prime}=\alpha^{2} e^{\alpha x}$ より, 特性方程式 $\alpha^{2}+2 \alpha+2=0$ を得る. 2 次方程式 $a x^{2}+2 b^{\prime} x+c=0$ の解の公式より

$$
x=\frac{-b^{\prime} \pm \sqrt{b^{\prime 2}-a c}}{a}
$$

であるから， $2$ 複素数解 $\alpha=-1 \pm \sqrt{1-2}=-1 \pm i$, （ $i$ は虚数単位 $\sqrt{-1}$ ) を得る. ここで $e^{(-1+i) x}$ と $e^{(-1-i) x}$ が互いに独立な解であるから, 求める一般解は, $c_{1}, c_{2}$ を任意定数として

$$
\begin{aligned}
y & =c_{1} e^{(-1+i) x}+c_{2} e^{(-1-i) x} \\
& \Longleftarrow \text { オイラー (Euler)の公式より } \\
& =c_{1} e^{-x}(\cos x+i \sin x)+C_{2} e^{-x}(\cos x-i \sin x) \\
& =e^{-x}\left(K_{1} \cos x+K_{2} \sin x\right)
\end{aligned}
$$

を得る. ただし, $K_{1}, K_{2}$ は任意定数である. また $K_{1}:=c_{1}+c_{2}$, および $K_{2}:=\left(c_{1}-c_{2}\right) i$ と置いた。

(3) $y^{\prime \prime}+2 y^{\prime}+1=0$ で $y=e^{\alpha x}$ と置いて, 特性方程式 $\alpha^{2}+2 \alpha+1=0$ を得る. これを 解いて $(\alpha+1)^{2}=0$ より, 重解 $\alpha=-1$ をもつ. このときは, $y=e^{\alpha x}$ 以外に $y=x e^{\alpha x}$ も独立な解となる. ゆえに, 求める一般解は

$$
y=c_{1} e^{-x}+c_{2} x e^{-x}=\left(c_{1}+c_{2} x\right) e^{-x}
$$

で与えられる. ただし, $c_{1}, c_{2}$ は任意定数である.

\bigskip
\bigskip
\bigskip
\bigskip
\bigskip
\bigskip
\bigskip
\bigskip
\bigskip
\bigskip
\bigskip
\bigskip
\bigskip
\bigskip

\begin{itembox}[l]{$[8]$ 次の微分方程式の一般解を求めよ.}
\begin{equation*}
y^{\prime \prime}-2 y^{\prime}+y=\sin x \tag{11}
\end{equation*}  
\end{itembox}

\bigskip

\begin{answer*}
\end{answer*}

まず斉次方程式の一般解を求める. $y=e^{\alpha x}$ として, 特性方程式 $\alpha^{2}-2 \alpha+1=0$ を得る. これを解いて $(\alpha-1)^{2}=0$ であるから, 重解 $\alpha=1$ をもつ. このとき, $e^{x}$ と $x e^{x}$ が独立 な 2 解を与える. ゆえに, 斉次方程式 $y^{\prime \prime}-2 y^{\prime}+y=0$ の一般解は

$$
y=\left(c_{1}+c_{2} x\right) e^{x}
$$

で与えられる. ただし $, c_{1}, c_{2}$ は任意定数である.

次に与式 $(11)$ の特別解 $y$ を求める. そのため,


\begin{equation*}
y=A \sin x+B \cos x \tag{*1}
\end{equation*}


と置いて, $y^{\prime}=A \cos x-B \sin x, y^{\prime \prime}=-A \sin x-B \cos x$ となる. これらを元の方程式 に代入して

$$
\begin{aligned}
y^{\prime \prime}-2 y^{\prime}+y & =(-A \sin x-B \cos x)-2 A \cos x+2 B \sin x+(A \sin x+B \cos x) \\
& =2 B \sin x-2 A \cos x=\sin x
\end{aligned}
$$

であるはずだから, 直ちに $A=0$ かつ $B=\frac{1}{2}$ となる.したがって, 非斉次方程式の特殊解は


\begin{equation*}
y=\frac{1}{2} \cos x \tag{*2}
\end{equation*}


である.ここで求める非斉次方程式 (11) $y^{\prime \prime}-2 y^{\prime}+y=\sin x$ の一般解は線形方程式の一般論から

$$
\text { 「特殊解」 }+ \text { 「斉次方程式の一般解」 }
$$

の形で与えられる. ゆえに $[8]$ の非斉次方程式 (11) の解は


\begin{equation*}
y(x)=\frac{1}{2} \cos x+\left(c_{1}+c_{2} x\right) e^{x} \tag{*3}
\end{equation*}


で与えられる.

\bigskip
\bigskip

\begin{itembox}[l]{[9]}
$\left\{f_{n}(x)\right\}_{n}$ を区間 $I=[a, b]$ 上で定義された連続関数列とする. この $\left\{f_{n}(x)\right\}_{n}$ がある 関数 $f(x)$ に $I$ 上で一様収束するとき,


\begin{equation*}
\lim _{n \rightarrow \infty} \int_{a}^{b} f_{n}(x) d x=\int_{a}^{b} f(x) d x \tag{12}
\end{equation*}


が成り立つことを示せ.  
\end{itembox}


\bigskip

\begin{answer*}
\end{answer*}

任意の$\varepsilon>0$ を取る.このとき, $f_n$ は$f$に一様収束するので

$\exists N \in \mathbb{N}$ s.t. $\forall x \in I, \ n \ge \mathbb{N}$ $\Rightarrow\left|f_{n}(x)-f(x)\right|< \displaystyle \frac{\varepsilon}{b-a}$
となる.
$f_n$は連続関数なのでRiemann積分可能

であり, 連続関数の一様収束先も連続関数なので$f$も連続関数となりRiemann積分可能.

$N \in \mathbb{N}$ に対して, \ $n \ge \mathbb{N}$ ならば

$$
\begin{aligned}
& \left|\int_{a}^{b} f_{n}(x) d x-\int_{a}^{b} f(x) d x\right|
=\left|\int_{a}^{b}\left(f_{n}(x)-f(x)\right) d x\right| 
\le \int_{a}^{b}\left|f_{n}(x)-f(x)\right| d x \\
& < \int_{a}^{b} \frac{\varepsilon}{b-a} d x
=\varepsilon
\end{aligned}
$$

よって問の等式が成立する.

\bigskip
\bigskip

\begin{itembox}[l]{[10]}
$I=[a, b] \subset \mathbb{R}$ とする. $\mathbb{R}^{2}$ 内の有界閉集合 $D=I \times E$ 上で定義された関数 $f(x, y)$ は連続であり, その導関数 $\frac{\partial f}{\partial y}$ も $D$ で連続であるとする.このとき, 任意の $\left(x_{0}, c\right) \in D$ において, 初期値問題

\[
\left\{\begin{array}{l}
\frac{d y}{d x}=f(x, y)  \tag{13}\\
y\left(x_{0}\right)=c
\end{array}\right.
\]

をみたす解 $y \equiv y(x)$ が $x=x_{0}$ の近傍 $U$ でただ 1 つ存在することを示せ.  
\end{itembox}

\bigskip

\begin{answer*}
\end{answer*}

$D$ は $\mathbb{R}^{2}$ 内のコンパクト集合で, $f \equiv f(x, y)$ の $D$ 上連続性より

$$
\exists M>0: \quad|f(x, y)| \leqslant M \quad(\forall(x, y) \in D)
$$

が成り立つ. すなわち, $f(x, y)$ は $D$ 上, $M$-有界である.

仮定より, $f \in C(D)$ は $D$ で連続な導関数 $\frac{\partial f}{\partial y}$ をもつから, $f_{x} \equiv f(x, \cdot) \in C^{1}(E)$ である.

したがって, $E$ 内の適当な任意の区間 $[\xi, \eta]$ において, 適当な点 $\exists \zeta(\xi<\zeta<\eta)$ がとれて

$$
\begin{aligned}
& \frac{f(x, \eta)-f(x, \xi)}{\eta-\xi}=\frac{\partial f}{\partial y}(\zeta) \quad(\forall x \in I) \\
& \Longleftrightarrow \quad f(x, \eta)-f(x, \xi)=\frac{\partial f}{\partial y}(\zeta) \cdot(\eta-\xi)
\end{aligned}
$$

仮定から $f(x, \cdot) \in C^{1}(E)$ であるから

$$
\begin{aligned}
& \exists K_{x}:=\max _{y \in[\xi, \eta]}\left|\frac{\partial f}{\partial y}(x, y)\right|, \quad(\forall x \in I) \\
& \text { i.e., }|f(x, \eta)-f(x, \xi)| \leqslant\left|\frac{\partial f}{\partial y}(\zeta)\right| \cdot|\eta-\xi| \quad(\forall x \in I) \\
& \leqslant K \cdot|\eta-\xi|
\end{aligned}
$$

が成り立つ.これは関数 $f(x, y)$ の変数 $y$ に関する一様リプシッツ(Lipschitz) 条件に他ならない.

したがって, 初期値問題をみたす解 $y \equiv y(x)$ が $x=x_{0}$ の近傍 $U$ でただ 1 つ存在する


\end{document}
